\section{引言}
\label{sec:introduction}

短期风电功率预测需要在预测起点可获得的信息基础上，估计未来一段时间内的发电轨迹。历史功率与数值天气预报（numerical weather prediction，NWP）是两类重要输入，但二者提供的信息具有不同含义：历史功率直接记录已经实现的发电水平及其变化，未来 NWP 则描述尚未发生的气象条件。如何组织这两类信息，使天气预测能够有效修正由历史功率支持的未来判断，是本文关注的核心问题。

这一问题在天气与功率具有不同时间分辨率时尤为突出。本文利用逐小时 NWP，预测未来六小时内每十五分钟的风场容量因子。线性插值能够将天气变量对齐至预测时刻，却不能单独回答这些变量应如何影响功率预测。对于相同的未来天气，不同历史发电轨迹所对应的预测修正可能不同。因此，修正网络除了需要天气特征，还需要明确知道它正在修正怎样一条未来功率轨迹。

本文将这一参考显式表示在功率空间中。首先，历史预测器生成未来多步基准轨迹；随后，另一网络结合未来天气、该基准轨迹以及历史表示，输出加性的功率修正。基准轨迹由历史信息及可用上下文计算，并未引入独立的新观测。它提供经预测任务监督、与未来目标对齐的参考表示，使天气分支能够围绕已有预测形成修正。由此，首要研究问题是：\emph{相较于只提供天气特征或历史隐表示，显式提供待修正的历史预测轨迹，能否提高预测准确率？}

显式轨迹与历史隐表示也不应简单视为等价输入。前者描述历史模型对各未来时刻功率水平的判断；后者保留已经实现的历史序列信息，其中一部分未必能从未来预测轨迹中恢复。本文使用独立的历史 Transformer 产生完整历史记忆，并以其最后一个上下文化 token 作为条件摘要。在每个预测时刻，该摘要与基准水平、相邻基准差分及提前量共同构成条件表示；条件与对齐后的天气特征进入普通非线性适配器。完整历史记忆则作为解码器的另一输入，参与未来功率修正的生成。

本文将上述框架称为\emph{历史与预测轨迹条件化的天气修正}（History/Trajectory-Conditioned Weather Correction）。训练先估计历史预测器，再固定其参数、训练修正网络。冻结参考使修正学习具有稳定的预测对象，但不意味着基准是在“没有天气影响”的物理世界中生成的反事实。历史功率已包含过去气象与运行条件的影响，修正项也可能同时承担天气相关预测与参考校准作用。本文据此区分预测收益、计算依赖与物理因果解释，不引入未经验证的动力学方程。

现有实验中，显式参考条件是证据较强的设计；额外历史摘要入口的收益较小，并随读出方式变化。最后一个上下文化 token 提供了无需新增读出参数的选择，而显式乘法门控及大型摘要网络尚未显示足以优先采用的收益。本文因而将重点放在信息组织和条件入口，而非特定交互算子的复杂度。

主要贡献包括：
\begin{itemize}
 \item 提出以历史预测轨迹为显式参考的天气条件修正框架，使基准同时进入修正条件和最终输出，形成具有固定预测参照的多步功率预测结构。
 \item 构建同时保留参考轨迹、历史上下文摘要与完整历史记忆的紧凑实现，通过直接天气路径和条件调整路径生成未来表示，联合预测各提前量的功率修正。
 \item 通过分协议对比、条件入口消融、摘要读出比较及参考路径诊断，量化不同表示的增量，对乘法交互、轨迹差分及历史摘要的作用给出与证据强度相符的解释。
\end{itemize}

本文以 56 个风场共享参数的容量因子预测为主要任务。区域功率可由风场预测按装机容量换算后求和，但风场误差与区域聚合误差对应不同评价目标。以下介绍模型方法及已有实验结果；相关工作综述暂不纳入本稿。
